%% ============================================================================
%% ch8_summary.tex -- Chapter 8: פרק הסיכום: מוכנים לשלב המימוש
%% budget: body page 37 (1 page)
%% ============================================================================
\chapter{סיכום ומוכנות למימוש}
\label{ch:summary}

\begin{managerbox}
האם אפשר להתחיל במימוש, ומה עדיין צריך להיבדק לפני שמתחייבים לתוצאה?
\end{managerbox}

הדוח הראה, בשרשרת אחת, שהפרויקט מוכן לשלב המימוש (\en{Design Readiness}). \textbf{הצורך ממשי ומכומת:} שגיאת צבע ממוצעת של \num{$\Delta E_{00} = 6.58$}, פי \num{3.3} מסף ההבחנה, שגיאת היקף גדולה ממידה אחת, והדמיה מסחרית שאין מנגנון שמממן אותה (פרק~\ref{ch:problem}). \textbf{הבסיס המדעי קיים והפער מוגדר:} ארבעה תחומי ידע מפותחים, והנחה משותפת שלא נבדקה, שצבעי תמונת הקלט נכונים (פרק~\ref{ch:lit}). \textbf{הכדאיות הטכנו-כלכלית נבחנה:} בהנחות שטרם אומתו, \en{SOM} של כ-\num{\$45.0M} משני ערוצים נפרדים, ו-\en{LTV/CAC} בחציון של \num{6.1}, עם סיכוי של \num{78\%} לעבור את סף \num{3.8} (פרק~\ref{ch:market}). \textbf{הארכיטקטורה נבחרה בשיקול גלוי:} החלופה ההיברידית קיבלה \num{4.30} מול \num{3.40} ו-\num{2.85}, ונשארה מובילה בכל בדיקות הרגישות למשקלים ולציונים (פרק~\ref{ch:design}). \textbf{התוכנית והסיכונים ממופים:} \num{12} משימות, נתיב קריטי, שתי אבני דרך, ושישה סיכונים קריטיים שציונם השיורי ירד לאזור הבינוני או הנמוך (פרקים~\ref{ch:plan} ו-\ref{ch:risk}).

טבלה~\ref{tab:evidence} מפרידה בין מה שכבר קיים לבין מה שעדיין הנחה, כדי שאף קורא לא יטעה בין סימולציה למדידה.

\begin{reporttable*}[\footnotesize]{מצב הראיות: מה קיים היום ומה עדיין פתוח}{tab:evidence}{P{2.9cm} R R}
\headrow \hc{נושא} & \hc{מה קיים היום} & \hc{מה עדיין פתוח} \\ \midrule
כיול צבע & סימולציה ספקטרלית וניסוי מדומה & ניסוי של \num{120} משתתפות; סיכוי של \num{12\%} בלבד שהחציון ביעד \\
מידות גוף & סימולציה של שילוב בייסיאני & מדידות סרט אמיתיות (\en{REQ-BOD-03}) \\
הדמיה & מאמר דגל עם קוד ומשקלים פתוחים & שחזור וזיקוק לשמונה צעדים (\en{TSK-3.1}) \\
כלכלה ושוק & מודל שכבות על הנחות צוות & המרה, נטישה ועלות גיוס מנתוני שותף \\
שותף ומשתמשות & שותף מדומה; אות ביקוש אורגני ברשתות, בלי נתוני משתמשות & הסכם פיילוט, אישור אתיקה והמרה מצפייה להתקנה \\
\end{reporttable*}

\textbf{החלטות פתוחות, לפי סדר עדיפות.} (\num{1}) הרחבת ניסוי האימות מארבע לשמונה משפחות תאורה, המדידה בעלת ערך המידע הגבוה ביותר (סעיף~\ref{sec:posterior}); אם החציון יישאר מעל \num{1.8}: עוגן שלישי או ניסוח מחדש של \en{REQ-VIS-01}, באבן דרך א'. (\num{2}) אם הזיקוק לא יעמוד ב-\num{4.5} שניות: רזולוציה נמוכה או הדמיה בפרימיום בלבד. (\num{3}) ניסוי ההמרה אצל השותף, לפני כל מדידת עלות (סעיף~\ref{sec:unit}). (\num{4}) שותף ממשי במקום המדומה. (\num{5}) מסגרת משפטית: הסכמת קטינות ואחריות להמלצה שגויה. \textbf{מחוץ לתחום:} אבחון רפואי, השקה מסחרית והצינור המרחבי (סעיף~\ref{sec:roadmap}). שלוש הפעולות שיוצאות לפועל ביום אישור הדוח מוצגות באיור~\ref{fig:actions}.

{\renewcommand{\needspace}[1]{}%
\begin{reportfigure*}{שלוש הפעולות המיידיות אחרי אישור הדוח}{fig:actions}
\begin{tikzpicture}[node distance=0.5cm,
  b/.style={draw=#1,fill=#1!10,rounded corners=3pt,minimum height=1.05cm,text width=2.75cm,
            align=center,font=\small,line width=0.9pt}]
  \node[b=brandsecondary] (a0) {\heb{הדוח אושר}};
  \node[b=brandaccent,left=of a0] (a1) {\heb{בקשת אתיקה והרשאות נתונים}};
  \node[b=brandaccent,left=of a1] (a2) {\heb{הקפאת תצורה וענפי פיתוח במאגר}};
  \node[b=brandaccent,left=of a2] (a3) {\heb{פגישת התנעה עם המנחה}};
  \foreach \a/\b in {a0/a1,a1/a2,a2/a3} \draw[-Stealth,thick,brandprimary] (\a) -- (\b);
\end{tikzpicture}
\end{reportfigure*}}

